\section{Account of the Events}\label{sec:relato}
This section reconstructs, in chronological order, what public agencies, researchers and the press reported, with each fact attributed to its source. Times are local (UTC$-$3) and approximate. Where sources differ, all accounts are presented without resolving the differences.

\subsection*{Background}
\paragraph{The Catchment and Alluvial Fan.} According to the analysis by Centro de Ciencia del Clima y la Resiliencia (CR2), Quebrada Honda drains a small catchment in the Andean foothills, covering about \SI{10.6}{km^2} and ranging from \SI{929}{m} to \SI{2355}{m} in elevation, which discharges at Santiago's urban boundary \citep{cr2_2026}. According to La Tercera, at its foot lies the San Carlos de Apoquindo neighborhood, where schools, clinics and universities are concentrated \citep{latercera2026zonacero}. The dean of Architecture at Universidad del Desarrollo, quoted by El Desconcierto, explained that several hillside ravines converge into a single channel and that ``the ravine is routed into a conduit at Bulnes Correa, and the problem with enclosed culverts is that they have limited capacity'' \citep{eldesconcierto2026}. As he told The Clinic, that channel runs through a park alongside General Blanche and, once enclosed, continues beneath Quebrada Honda street \citep{theclinic2026allard}.

\paragraph{The Fire.} According to CR2, a fire began in the catchment on 29 December 2025 \citep{cr2_2026}. The municipality reported that it started in the La Plaza Sur sector, affected a preliminary area of \SI{1012}{ha} within the municipality, and was declared under control by Conaf on the afternoon of 1 January \citep{lascondes2026}. Area figures vary: ``about 1,000 hectares'' according to CR2, which adds that the fire affected virtually the entire catchment \citep{cr2_2026}; ``approximately one thousand hectares'' over four days according to T13 \citep{t13_2026mallas}; and ``more than 800 hectares'', with the fire within \SI{200}{m} of urban areas, according to Interferencia \citep{interferencia2026}. According to the calculation in this study (Section~\ref{sec:res-incendio}), \SI{8.2}{km^2} of the \SI{10.1}{km^2} of contributing catchments show a burn signal, mostly of low or moderate-low severity.

\paragraph{The Event of 31 January.} According to the Sernageomin report, that day, flows were triggered in the La Honda, Grande, Los Almendros and Apoquindo ravines \citep{sernageomin2026reterm}. According to T13, \SI{75}{mm} fell in one hour and \SI{20}{min}, and the event affected eight urban hectares, almost \SI{4}{km} of streets and 11 homes \citep{t13_2026mallas}. Interferencia, however, mentions ``just 20 minutes of intense rain'' \citep{interferencia2026}. El Mostrador puts the number of damaged homes at 22, according to initial information from Senapred \citep{elmostrador2026}. The mayor of Las Condes said at the time, according to La Tercera, that the Francisco Bulnes Correa culvert ``worked very well'', although some locations were affected, including Mirasol and Carlos Peña streets \citep{latercera2026enero}.

\paragraph{Sernageomin's Warning.} According to the report itself, following an overflight on 6 February, Servicio Nacional de Geología y Minería prepared the Reporte Técnico de Remoción en Masa Metropolitana 02/2026 on 11 February \citep{sernageomin2026reterm}. According to Chilevisión, the report placed the area in a zone of ``high hazard'' from debris flows and flooding, identified flows in the La Honda, Grande and Apoquindo ravines, and noted ``deficiencies in mitigation and adaptation works'' \citep{chilevision2026}. The Sernageomin report concluded that the January event ``does not constitute an isolated occurrence, but a reactivation of processes in an area highly prone'' to these phenomena \citep{sernageomin2026reterm}. According to Chilevisión, among other measures, the report recommended evacuation plans and periodic channel cleaning \citep{chilevision2026}.

\paragraph{Mitigation Works.} According to Interferencia, on 26 March, the mayor asked the presidential delegate to install flexible mesh barriers; he acknowledged that it would be difficult to have them before May \citep{interferencia2026}. According to the same outlet, on 7 August an agreement was announced for five barriers with \$754 million from Ministerio de Obras Públicas (MOP) and a 90-day timeframe \citep{interferencia2026}. T13 reports a different figure: a directly awarded contract for \$652 million, signed in September, also with a 90-day timeframe \citep{t13_2026mallas}. According to Interferencia, during the emergency the presidential delegate confirmed that they were still not ready \citep{interferencia2026}.

\subsection*{The Rainfall}
According to CR2, a category 2--3 atmospheric river reached central Chile on 6 and 7 October. It brought \SI{72}{mm} to Quinta Normal and about \SI{100}{mm} to the Santiago foothills \citep{cr2_2026}. On the 8th, a stationary cutoff low persisted at an altitude of about \SI{5}{km}; according to CR2, together with moisture from the atmospheric river, it generated convective cells with sudden darkening, lightning, hail and intense rainfall over periods of \SIrange{30}{60}{min} \citep{cr2_2026}. CR2 notes that there appear to be no rain gauges in the catchment; at Río Molina station, about \SI{7}{km} away, \SI{42.8}{mm} fell that day, including \SI{31.6}{mm} between 15:00 and 19:00, with a maximum of \SI{10.9}{mm} in \SI{30}{min}. According to that analysis, the \SI{0}{\celsius} isotherm was at about \SI{2800}{m}, so the entire catchment received rain rather than snow \citep{cr2_2026}. According to CR2, lightning reached northeastern Santiago and Las Condes around 16:30, and by about 17:00 it was concentrated mainly over Cajón del Maipo \citep{cr2_2026}. At Tobalaba, Dirección Meteorológica de Chile recorded \SI{34.8}{mm} on the 8th and \SI{13.8}{mm} between 15:00 and 21:00 \citep{dmc2026}.

Other figures circulated without an identified measurement. The mayor spoke of \SI{144}{mm} of accumulated rainfall, \SI{20}{mm} of it that day, with saturated soil \citep{interferencia2026}. According to the version of Wikipedia consulted, the undersecretary of the Interior mentioned about \SI{15}{mm} in a very short interval \citep{wikipedia2026}. Press searches logged on 9 October also returned figures of \SIrange{16}{20}{mm} in about one hour, with no station identified. The figure of \SI{27}{mm} in \SI{20}{min} was not found in the press consulted and is unverified. The \SI{25}{mm} figure that also circulated appears to correspond to a forecast for 11--13 October broadcast by Canal 9, rather than rainfall on the 8th (search logged on 12 October).

\subsection*{The Debris Flow}
Around 17:00, according to preliminary information from Carabineros, a ``large torrent of water along General Blanche street'' was sighted, which that information attributed to overflow from the channel on Francisco Bulnes Correa street; the flow descended westward \citep{theclinic2026frontal,eldesconcierto2026}. Sources place the origin at different locations. La Tercera reported that ``it all began when a ravine became active at General Blanche and San Carlos de Apoquindo'', with the flow advancing southwestward \citep{latercera2026}. According to La Tercera, using a drone, Carabineros located ``ground zero'' at Mons.\ Álvaro del Portillo and Av.\ La Plaza \citep{latercera2026zonacero}. According to Chilevisión, the authorities attributed the event to the activation of ``Quebrada Las Ondas'' \citep{chilevision2026mapa}, whereas Chilevisión, drawing on the Sernageomin report, and CR2 refer to the La Honda or Honda ravine \citep{chilevision2026,cr2_2026}.

The vice-rector for Communications at Universidad de los Andes described to Cooperativa an overflow from the ravine on the north side of the campus, bordering General Blanche, ``down which this entire mass of mud fell'' \citep{cooperativa2026}. According to El Desconcierto, the flow swept away cars and a RED bus, whose passengers had to climb onto the roof \citep{eldesconcierto2026}. La Tercera and T13 reported that the bus overturned \citep{latercera2026,t13_2026vehiculos}; Cadena 3 reported that a public transport vehicle managed to stop \citep{cadena3_2026}. According to the municipality, mud descended along Camino El Alba, General Blanche and Quebrada Honda \citep{latercera2026,chilevision2026mapa}. According to El Dínamo, Senapred identified General Blanche, Quebrada Honda and Av.\ La Plaza Sur as the most affected streets \citep{eldinamo2026}. According to T13, water also accumulated on Av.\ Padre Hurtado \citep{t13_2026vehiculos}.

At Clínica UC Christus, Cadena 3 reported about \SI{1}{m} of mud in the emergency department \citep{cadena3_2026}. According to El Desconcierto, services in that department and consultations were suspended \citep{eldesconcierto2026}. According to BioBioChile, at the Duoc UC campus, students and staff were trapped in an auditorium \citep{biobio2026duoc}. According to Cooperativa, the Universidad de los Andes community remained on campus on the recommendation of Bomberos \citep{cooperativa2026}. A video circulated on Facebook is titled ``Debris flow entered the Aula Magna of Universidad de Los Andes'' \citep{quintaconectado2026}, and Meganoticias published images of the flooded campus at 18:16 \citep{meganoticias2026uandes}. Interferencia also reported water and mud at Hospital de la FACh \citep{interferencia2026}.

Different accounts of the drainage system were reported, all preliminary. According to Chilevisión, the presidential delegate presumed ``a blockage, a plugging of the underground channel'' at Álvaro del Portillo and Av.\ La Plaza \citep{chilevision2026}. According to Interferencia, in another statement the delegate pointed to the obstruction of two conduits on General Blanche \citep{interferencia2026}. According to El Desconcierto, the enclosed channel descending from Av.\ La Plaza reaches a settling chamber at Francisco Bulnes Correa and, ``with that system collapsed, the street served as the new channel'' \citep{eldesconcierto2026}. Interferencia emphasized that this was a hypothesis without technical confirmation \citep{interferencia2026}.

\subsection*{Consequences and Response}
\paragraph{The Afternoon and Night of the 8th.} According to Interferencia, Senapred issued its first public notice at 17:41, with a video of its regional director; it did not send an evacuation message to mobile phones \citep{interferencia2026}. The red alert for ten municipalities took effect at 18:20 according to Interferencia \citep{interferencia2026}; La Tercera places it at 18:30 \citep{latercera2026lluvias}. According to La Tercera, at 18:15 the president of the Republic wrote that they were facing ``an emergency with a very rapid onset'' \citep{latercera2026lluvias}. According to Interferencia, that night the president ordered a disaster-area declaration following a regional risk management committee meeting \citep{interferencia2026}; according to the version of Wikipedia consulted, the measure was formalized on 9 October \citep{wikipedia2026}. The presidential delegate reported, according to Interferencia and El Desconcierto, six injured people, including a one-year-old child with a fracture, and no deaths \citep{interferencia2026,eldesconcierto2026}. According to Chilevisión, the minister of the Interior announced an exclusion zone following aerial reconnaissance \citep{chilevision2026mapa}. After 23:00, the mayor told Cooperativa that ``the hill gave way almost completely'' \citep{cooperativa2026alcaldesa}.

\paragraph{The Perimeter.} According to The Clinic, the municipality's first safety perimeter extended from Camino El Alba to Quebrada Honda, with Vital Apoquindo and the mountain boundary to the east, subject to adjustments \citep{theclinic2026frontal}. According to La Tercera, the polygon reported later was bounded by Camino El Alba (north), Av.\ La Plaza (east), Quebrada Honda (south) and Av.\ Padre Hurtado (west), in the Pueblito Los Dominicos sector \citep{latercera2026}.

\paragraph{Area and Homes.} Figures differ by agency and date:
\begin{itemize}
\item Area: about \SI{200}{ha} according to the assessment for the night of the 8th reproduced in the consulted Wikipedia version \citep{wikipedia2026}; \SI{253}{ha} and 3,981 property tax assessment identifiers according to the real estate platform Inciti \citep{chicureohoy2026}; \SI{288.5}{ha}, 2,266 properties and 368 street segments in the official IDE Chile viewer \citep{bbnn2026visor}.
\item Municipality (9 October): 2,015 affected homes, 55 uninhabitable and 1,960 with minor damage, according to CNN Chile and Chicureo Hoy \citep{cnnchile2026balance,chicureohoy2026}. That same day, the presidential delegate spoke of 1,960 affected homes, 55 of them with ``serious habitability problems'', according to Publimetro \citep{publimetro2026}.
\item Senapred: 75 homes with major damage and 2,286 with minor damage on 9 October, according to La Tribuna \citep{latribuna2026}; later, according to La Tercera, 56 destroyed, 40 with major damage, 1,960 with minor damage and 384 disaster-affected people \citep{latercera2026senapred}.
\item Bienes Nacionales viewer: 22 homes with major damage, 1,960 with minor damage and 51 affected facilities \citep{bbnn2026visor}.
\end{itemize}
The municipality put the number of people evacuated or rescued at 1,200, according to CNN Chile \citep{cnnchile2026balance}.

\paragraph{Clínica Universidad de los Andes.} Accounts differ. Meganoticias reported that mud reached the emergency department \citep{meganoticias2026urgencia}, and BioBioChile circulated images of flooded rooms and corridors attributed to this clinic and to Clínica UC Christus \citep{biobio2026clinicas}. The clinic's director of Communications stated to Radio La Señal that ``all the impact was outside'' and that only ``a little water in a waiting room'' had entered, no more than one centimeter deep \citep{radiolasenal2026}. The municipal assessment recorded ``a minor ingress of water'', with the emergency department available, according to CNN Chile \citep{cnnchile2026balance}.

\paragraph{The Mesh Barriers.} According to T13, on 9 October the minister of Public Works said that there had been ``installations of flexible mesh barriers''. The mayor then clarified that access routes still had to be prepared and that installation ``was starting tomorrow''; the minister later acknowledged that the works were at a preliminary stage \citep{t13_2026mallas}. According to Chicureo Hoy, Ministerio Público opened an investigation into possible criminal liability \citep{chicureohoy2026}.

\subsection*{Questions Raised by the Events}
The sources cited in this account disagree on where the flow entered the city: Álvaro del Portillo and La Plaza, General Blanche and San Carlos de Apoquindo, or the Francisco Bulnes Correa channel. Nor is it established whether the mud came from a single channel or several ravines, as a Las Condes councilor suggested to CNN Chile \citep{cnnchile2026haddad}. No measurements were located in the sources consulted of rainfall within the catchment or of the volume mobilized, which the presidential delegate described as ``thousands of cubic meters'' \citep{cooperativa2026alcaldesa}. Published areas and counts differ, and none report thicknesses or high-water marks. These questions motivate the hypotheses in Section~\ref{sec:hipotesis}.
